\chapter{Objetivos del TFM}
% BORRADOR para editar a tu voz. Recomendado: 1-3 páginas.

En este capítulo se establece el objetivo general del trabajo y los objetivos específicos en que
se desglosa. Los objetivos se redactan en infinitivo y se corresponden con metas concretas
alcanzadas al término del proyecto: uno asociado al marco teórico, otro al desarrollo
metodológico y los restantes a la validación empírica.

\section{Evolución del alcance del proyecto}
El planteamiento inicial situaba la explicabilidad visual —la morfología del mapa de calor de
atención— como la señal candidata para estimar la fiabilidad de un informe. La evidencia recabada
en las primeras fases del trabajo (Capítulo~5, PI2) mostró que dicha morfología no discrimina entre
aciertos y errores, lo que exigió una reformulación del alcance: el foco se desplazó a la
incertidumbre de la decisión como señal alternativa, y el diseño experimental se amplió a una
validación cruzada \(2\times2\) (dos generadores por dos dominios) para descartar que el resultado
negativo de la explicabilidad fuese un artefacto de un único modelo o población. El propio orden de trabajo responde a una decisión deliberada: se optó por construir primero un
\textbf{generador propio}, más pequeño y con destilación de conocimiento, precisamente para entender
\emph{desde dentro} el funcionamiento de estos sistemas y las limitaciones específicas de un modelo
de menor capacidad, antes de dar el salto a un modelo del estado del arte ya entrenado y tratado como
caja cerrada. Esa fase reveló, entre otras lecciones, que las métricas de solapamiento léxico como
ROUGE son inadecuadas para valorar el equilibrio entre acierto y alucinación (Sección~5.2), lo que
motivó abandonarlas en favor de la corrección clínica. Junto con la combinación de cómputo en
hardware de consumo y el acceso restringido a los grandes corpus clínicos (Sección~4.1), estas
lecciones llevaron a elevar el análisis de explicabilidad a un modelo del estado del arte —LLaVA-Rad,
y después MAIRA-2—, ejecutados localmente mediante cuantización en lugar de entrenados desde cero.
Los objetivos que siguen recogen ya el alcance resultante de esta reformulación, no el planteamiento
original.

\section{Objetivo general}
Evaluar la fidelidad y el anclaje clínico de las explicaciones visuales en la generación
automática de informes de radiografía de tórax mediante arquitecturas encoder--decoder
multimodales, y determinar en qué medida dichas explicaciones —o la incertidumbre del propio
modelo— permiten estimar la fiabilidad de cada informe.

\section{Objetivos específicos}
\begin{enumerate}
  \item \textbf{Analizar el estado del arte} en generación automática de informes radiológicos,
  en explicabilidad de modelos transformer y en detección de incertidumbre y alucinaciones,
  identificando el hueco que este trabajo aborda.
  \item \textbf{Diseñar e implementar} un sistema de generación de informes encoder--decoder
  multimodal, incorporando destilación de conocimiento desde modelos preentrenados libres como
  respuesta metodológica a la restricción de datos, junto con un marco reutilizable de
  explicabilidad.
  \item \textbf{Cuantificar de forma objetiva el anclaje clínico} de la atención de un modelo
  multimodal del estado del arte, contrastándola con anotaciones reales de localización
  (bounding boxes de NIH ChestX-ray8).
  \item \textbf{Determinar si la explicabilidad y la incertidumbre} del modelo permiten detectar
  sus errores (alucinaciones y omisiones), sentando las bases de una estimación de confianza por
  informe, agnóstica al modelo empleado.
\end{enumerate}
